\begin{abstract}
Underground Hydrogen Storage (UHS) in depleted reservoirs provides seasonal energy capacity, but transient pressure dynamics are computationally expensive to simulate. We present a data-driven surrogate framework on a 10,116-cell irregular PEBI mesh using five simulation cases. First, we expose an evaluation pitfall where direct pressure prediction yields deceptive accuracy ($R^2 > 0.99$), demonstrating that predicting pressure change ($\Delta P$) is a more rigorous target. Combining spatiotemporal features with a one-step pressure lag ($\Delta P_{t-1}$), a Random Forest model predicts local pressure changes with high accuracy ($R^2 = 0.9827$, RMSE = 0.5487 bar). To handle irregular mesh geometry, we propose a Spatio-Temporal Graph Neural Network (ST-GNN) mapping physical transmissibility to edge weights. In a 142-step autoregressive rollout, the optimized ST-GNN v3 (trained for 60 epochs using a multi-objective loss sweep) reduces the mean pressure RMSE by 23.0\% to 27.15 bar and final-step pressure RMSE to 31.66 bar, compared to the baseline v1 (35.26 bar mean, 49.25 bar final). However, key limitations remain, including the restricted five-case dataset under a single geological realization and long-horizon pressure drift peaking at 72.45 bar. These results show that while graph-based surrogates enable rapid screening, rollout stability and generalization remain primary challenges.
\end{abstract}
